% GENERATED FILE. Edit canonical lessons, then run npm run generate:books.
% canonical-source-hash: fnv1a64:551115cfc744ad50
% canonical-lessons: SA-C545-anukramah, SA-C545-yathakramam, SA-C545-rte, SA-C545-adau, SA-C545-alocayati

\chapter{In the Right Order}
\label{ch:sa-in-the-right-order}
\hlchaptermodality{545}

\begin{chapteropening}
\textbf{By the end of this chapter:} \emph{I can talk about order, steps and exceptions.}

Complete the last lesson of chapter 545: I can talk about order, steps and exceptions.
\end{chapteropening}

\section[\texorpdfstring{\sk{अनुक्रमः} (anukramaḥ)}{अनुक्रमः (anukramaḥ)}]{\sk{अनुक्रमः} (anukramaḥ) — a sequence}

\label{lesson:SA-C545-anukramah}



\begin{quote}
Before the new one: say the Sanskrit for separately, then the Sanskrit for at present.
\end{quote}

\subsection*{You'll want to know: \sk{अनुक्रमः}}
\textbf{\sk{अनुक्रमः}} — \emph{anukramaḥ} — \textquotedblleft{}a sequence\textquotedblright{}.

\begin{quote}
\sk{यदि} \sk{अनुक्रमः} \sk{सम्यक्} \sk{अस्ति}, \sk{तर्हि} \sk{कार्यं} \sk{सुकरं} \sk{भवति।}
\end{quote}

\begin{quote}
\emph{yadi anukramaḥ samyak asti, tarhi kāryaṃ sukaraṃ bhavati.}
\end{quote}

\begin{quote}
If the order is right, the work is easy.
\end{quote}

A masculine noun.

\begin{cousinweb}
\textbf{\sk{अनुक्रमः}} (\emph{anukramaḥ}) is built on \emph{anu}, \textquotedblleft{}after\textquotedblright{}, and \emph{krama}, \textquotedblleft{}a step\textquotedblright{}, from the root √kram, \textquotedblleft{}to step\textquotedblright{}; MW gives \emph{anukrama} as succession, arrangement, order (MW).
\end{cousinweb}

\begin{grammarlens}[title={What you've built}]
One more word for saying what happens if.
\end{grammarlens}

\subsection*{Your turn}
\begin{itemize}
\raggedright
  \item \emph{Say it:} \emph{anukramaḥ}
  \item \emph{Say it:} \emph{anukramaḥ}, once more
  \item \emph{Say it:} \emph{yadi anukramaḥ samyak asti, tarhi kāryaṃ sukaraṃ bhavati.} — If the order is right, the work is easy.
  \item \emph{Recall:} say the Sanskrit for separately, then the Sanskrit for at present, then say \emph{anukramaḥ} again
\end{itemize}

\begin{tcolorbox}[breakable,colback=teal!4,colframe=teal!35!black,title={Before you move on}]
What does \sk{अनुक्रमः} mean? (\textquotedblleft{}A sequence\textquotedblright{}.) Say it once more.
\end{tcolorbox}
\section[\texorpdfstring{\sk{यथाक्रमम्} (yathākramam)}{यथाक्रमम् (yathākramam)}]{\sk{यथाक्रमम्} (yathākramam) — in due order}

\label{lesson:SA-C545-yathakramam}



\begin{quote}
Before the new one: say the Sanskrit for at present, then the Sanskrit for a sequence.
\end{quote}

\subsection*{You'll want to know: \sk{यथाक्रमम्}}
\textbf{\sk{यथाक्रमम्}} — \emph{yathākramam} — \textquotedblleft{}in due order\textquotedblright{}.

\begin{quote}
\sk{यदि} \sk{सर्वे} \sk{यथाक्रमम्} \sk{आगच्छन्ति}, \sk{तर्हि} \sk{शब्दः} \sk{न} \sk{भवति।}
\end{quote}

\begin{quote}
\emph{yadi sarve yathākramam āgacchanti, tarhi śabdaḥ na bhavati.}
\end{quote}

\begin{quote}
If everyone comes in turn, there is no noise.
\end{quote}

\begin{cousinweb}
\textbf{\sk{यथाक्रमम्}} (\emph{yathākramam}) is built from \emph{yathā}, \textquotedblleft{}as\textquotedblright{}, and \emph{krama}, \textquotedblleft{}a step, an order\textquotedblright{}, from the root √kram; MW gives \emph{yathākramam} as according to order, in due succession (MW).
\end{cousinweb}

\begin{grammarlens}[title={What you've built}]
One more word for saying what happens if.
\end{grammarlens}

\subsection*{Your turn}
\begin{itemize}
\raggedright
  \item \emph{Say it:} \emph{yathākramam}
  \item \emph{Say it:} \emph{yathākramam}, once more
  \item \emph{Say it:} \emph{yadi sarve yathākramam āgacchanti, tarhi śabdaḥ na bhavati.} — If everyone comes in turn, there is no noise.
  \item \emph{Recall:} say the Sanskrit for at present, then the Sanskrit for a sequence, then say \emph{yathākramam} again
\end{itemize}

\begin{tcolorbox}[breakable,colback=teal!4,colframe=teal!35!black,title={Before you move on}]
What does \sk{यथाक्रमम्} mean? (\textquotedblleft{}In due order\textquotedblright{}.) Say it once more.
\end{tcolorbox}
\section[\texorpdfstring{\sk{ऋते} (ṛte)}{ऋते (ṛte)}]{\sk{ऋते} (ṛte) — except}

\label{lesson:SA-C545-rte}



\begin{quote}
Before the new one: say the Sanskrit for a sequence, then the Sanskrit for in due order.
\end{quote}

\subsection*{You'll want to know: \sk{ऋते}}
\textbf{\sk{ऋते}} — \emph{ṛte} — \textquotedblleft{}except\textquotedblright{}.

\begin{quote}
\sk{यदि} \sk{माम्} \sk{ऋते} \sk{सर्वे} \sk{आगच्छन्ति}, \sk{तर्हि} \sk{अहं} \sk{दुःखी} \sk{भविष्यामि।}
\end{quote}

\begin{quote}
\emph{yadi mām ṛte sarve āgacchanti, tarhi ahaṃ duḥkhī bhaviṣyāmi.}
\end{quote}

\begin{quote}
If everyone comes but me, I will be sad.
\end{quote}

\begin{cousinweb}
\textbf{\sk{ऋते}} (\emph{ṛte}) is the word MW gives for excepting, without, unless, and reports as perhaps a case form of the past participle of the root √ṛ, \textquotedblleft{}to go\textquotedblright{}, so its origin is uncertain (MW).
\end{cousinweb}

\begin{grammarlens}[title={What you've built}]
One more word for saying what happens if.
\end{grammarlens}

\subsection*{Your turn}
\begin{itemize}
\raggedright
  \item \emph{Say it:} \emph{ṛte}
  \item \emph{Say it:} \emph{ṛte}, once more
  \item \emph{Say it:} \emph{yadi mām ṛte sarve āgacchanti, tarhi ahaṃ duḥkhī bhaviṣyāmi.} — If everyone comes but me, I will be sad.
  \item \emph{Recall:} say the Sanskrit for a sequence, then the Sanskrit for in due order, then say \emph{ṛte} again
\end{itemize}

\begin{tcolorbox}[breakable,colback=teal!4,colframe=teal!35!black,title={Before you move on}]
What does \sk{ऋते} mean? (\textquotedblleft{}Except\textquotedblright{}.) Say it once more.
\end{tcolorbox}
\section[\texorpdfstring{\sk{आदौ} (ādau)}{आदौ (ādau)}]{\sk{आदौ} (ādau) — at first}

\label{lesson:SA-C545-adau}



\begin{quote}
Before the new one: say the Sanskrit for in due order, then the Sanskrit for except.
\end{quote}

\subsection*{You'll want to know: \sk{आदौ}}
\textbf{\sk{आदौ}} — \emph{ādau} — \textquotedblleft{}at first\textquotedblright{}.

\begin{quote}
\sk{यदि} \sk{आदौ} \sk{कठिनम्} \sk{अस्ति}, \sk{तर्हि} \sk{अनन्तरं} \sk{सुकरं} \sk{भविष्यति।}
\end{quote}

\begin{quote}
\emph{yadi ādau kaṭhinam asti, tarhi anantaraṃ sukaraṃ bhaviṣyati.}
\end{quote}

\begin{quote}
If it is hard at first, it will be easy later.
\end{quote}

\begin{cousinweb}
\textbf{\sk{आदौ}} (\emph{ādau}) is the locative of \emph{ādi}, \textquotedblleft{}a beginning\textquotedblright{}, which MW records as beginning, commencement; MW gives no further derivation, so its origin is uncertain (MW).
\end{cousinweb}

\begin{grammarlens}[title={What you've built}]
One more word for saying what happens if.
\end{grammarlens}

\subsection*{Your turn}
\begin{itemize}
\raggedright
  \item \emph{Say it:} \emph{ādau}
  \item \emph{Say it:} \emph{ādau}, once more
  \item \emph{Say it:} \emph{yadi ādau kaṭhinam asti, tarhi anantaraṃ sukaraṃ bhaviṣyati.} — If it is hard at first, it will be easy later.
  \item \emph{Recall:} say the Sanskrit for in due order, then the Sanskrit for except, then say \emph{ādau} again
\end{itemize}

\begin{tcolorbox}[breakable,colback=teal!4,colframe=teal!35!black,title={Before you move on}]
What does \sk{आदौ} mean? (\textquotedblleft{}At first\textquotedblright{}.) Say it once more.
\end{tcolorbox}
\section[\texorpdfstring{\sk{आलोचयति} (ālocayati)}{आलोचयति (ālocayati)}]{\sk{आलोचयति} (ālocayati) — reviews; considers}

\label{lesson:SA-C545-alocayati}



\begin{quote}
Before the new one: say the Sanskrit for except, then the Sanskrit for at first.
\end{quote}

\subsection*{You'll want to know: \sk{आलोचयति}}
\textbf{\sk{आलोचयति}} — \emph{ālocayati} — \textquotedblleft{}reviews; considers\textquotedblright{}.

\begin{quote}
\sk{यदि} \sk{त्वं} \sk{योजनाम्} \sk{आलोचयसि}, \sk{तर्हि} \sk{कार्यं} \sk{सम्यक्} \sk{भविष्यति।}
\end{quote}

\begin{quote}
\emph{yadi tvaṃ yojanām ālocayasi, tarhi kāryaṃ samyak bhaviṣyati.}
\end{quote}

\begin{quote}
If you go over the plan, the work will go well.
\end{quote}

Past (he …): \sk{सः} \sk{योजनाम्} \sk{आलोचितवान्}, \emph{saḥ yojanām ālocitavān}.

\begin{cousinweb}
\textbf{\sk{आलोचयति}} (\emph{ālocayati}) is built on \emph{ā} and the causative of the root √loc, \textquotedblleft{}to see\textquotedblright{}; MW gives \emph{ālocayati} as to behold, to consider, to reflect (MW).
\end{cousinweb}

\begin{grammarlens}[title={What you've built}]
One more word for saying what happens if.
\end{grammarlens}

\subsection*{Your turn}
\begin{itemize}
\raggedright
  \item \emph{Say it:} \emph{ālocayati}
  \item \emph{Say it:} \emph{ālocayati}, once more
  \item \emph{Say it:} \emph{yadi tvaṃ yojanām ālocayasi, tarhi kāryaṃ samyak bhaviṣyati.} — If you go over the plan, the work will go well.
  \item \emph{Recall:} say the Sanskrit for except, then the Sanskrit for at first, then say \emph{ālocayati} again
\end{itemize}

\begin{tcolorbox}[breakable,colback=teal!4,colframe=teal!35!black,title={Before you move on}]
What does \sk{आलोचयति} mean? (\textquotedblleft{}Reviews; considers\textquotedblright{}.) Say it once more.
\end{tcolorbox}
